\section{Minkowski-Weyl}

Various types of polyhedral objects (polytopes, polyhedra, polyhedral cones)
can be described in two complementary ways:
\begin{enumerate}
  \item A \emph{V-representation}, which describes it as the closure of a finite
  set of atoms (vertices) with respect to some notion
  (e.g., convex or conic combinations).
  \item A \emph{H-representation}, which describes it as the finite intersection
  of affine or linear halfspaces.
\end{enumerate}

The Minkowski--Weyl family of theorems says that these two types of
representations define more or less the same classes of objects,
sometimes under extra conditions such as boundedness.

The main objective of our work group was to formalize the H-representation of polyhedra
and proof the Minkowski-Weyl theorem which establishes its correspondence with the
V-representation, which was already defined for polytopes and cones.

One of the main obstacles we encountered in the formalization of both, the definition and
the result was the question on how to deal with infinite dimensions.
We are not aware of any existing convention on how the definition should generalize to this setting,
and none of the participants in the workshop expressed a strong interest in how this generalization
should be.
Two main possibilities were discussed:
\begin{enumerate}
  \item Formalize the classical definition of an H-polyhedron, as the finite intersection of
  halfspaces, and require finite dimension for some results, including the Minkowski-Weyl theorem.
  \item Admit the intersection with an arbitrary affine subspace.
  This makes it possible to formalize a version of the Minkowski-Weyl theorem that allows for
  an infinite amount on equalities on the H-representation and an infinite-dimensional lineality
  space on the V-representation, both of which collapse to the usual Minkowski-Weyl theorem on
  finite dimension.
\end{enumerate}

\subsection{Naming}
\label{subsec:mw:naming}
One of Mathlib's principles is to avoid having multiple definitions for the same concept.
The impact of a name choice affects not only usages of the definition, but also the way it is
referred to in theorem and lemma names.

An important question we considered in the workshop was whether there should be two separate names
for \texttt{VPolytope}/\texttt{HPolytope}, as well as for \texttt{VPolyhedron}/\texttt{HPolyhedron}.
%and \texttt{VCone}/\texttt{HCone}.
After several discussions, we reached a consensus on the choice
to define only \texttt{Polytope} and \texttt{Polyhedron} as the V-representation and
H-representation respectively, and perhaps define the alternatives \texttt{HPolytope} and
\texttt{VPolyhedron} only if it made a project's work easier.

This decision is motivated by the fact each representation is the most natural for each concept.
The H-representation of polytopes requires the notion of boundedness, either algebraic or
topological, while the V-representation of polyhedra requires atoms of two types, vertices and
rays.

Regardless of this choice, API design should provide convenient access to the V-representation
and H-representation of both definitions, which relies on the Minkowski--Weyl theorem we aim to
formalize.

On the topic of polyhedral cones, no such decision was reached because,
on the one hand, both definitions are equally natural, and on the other hand,
these notions are already defined in Mathlib as finitely generated (\texttt{PointedCone.FG}) and
dually finitely generated cones (\texttt{PointedCone.DualFG}).

\subsection{Hypotheses}
\label{subsec:mw:hypotheses}
Our work follows the previous efforts to formalize the Minkowski--Weyl theorem for cones,
which was already stated in terms of finitely generated cones and dually finitely generated
cones over fields.

These results use the \texttt{PointedCone.FG} and \texttt{PointedCone.DualFG} definitions,
which are defined for modules over a commutative ring.
\begin{lstlisting}
variable {R M N : Type*}
variable [CommRing R] [PartialOrder R] [IsOrderedRing R]
variable [AddCommGroup M] [Module R M]
variable [AddCommGroup N] [Module R N]
\end{lstlisting}

The first one, \texttt{FG}, is simply an abbreviation for \texttt{Submodule.FG},
in the sense that \texttt{PointedCone} is defined as a submodule over the nonnegative semiring.
\begin{lstlisting}
abbrev PointedCone (R E)
    [Semiring R] [PartialOrder R] [IsOrderedRing R] [AddCommMonoid E] [Module R E] :=
  Submodule {c : R // 0 ≤ c} E

def Submodule.FG (N : Submodule R M) : Prop :=
  ∃ S : Finset M, span R ↑S = N

abbrev FG (C : PointedCone R M) : Prop := Submodule.FG C
\end{lstlisting}

On the other hand, \texttt{DualFG} is defined using a bilinear pairing \texttt{p}
and the notion of dual cone.
The dual cone of a set $S$ with respect to a bilinear pairing $p$ is the cone
consisting of all points $y$ such that $0 \leq p(x, y)$ for all $x \in S$.

\begin{lstlisting}
/-- Bilinear map `(M × N) → R` -/
variable {p : M →ₗ[R] N →ₗ[R] R}

variable (p) in
def dual (s : Set M) : PointedCone R N where
  carrier := {y | ∀ ⦃x⦄, x ∈ s → 0 ≤ p x y}
  sorry -- proofs that the carrier is a PointedCone
\end{lstlisting}

This bilinear pairing serves as the generalization of an inner product.
A cone is dually finitely generated if it can be written as the dual cone of a finite set.

\begin{lstlisting}
variable (p) in
def DualFG (C : PointedCone R N) : Prop := ∃ s : Finset M, dual p s = C
\end{lstlisting}

In the sense of the Minkowski--Weyl theorem, a finitely generated cone is
a $V$-cone, as it is written as the conic closure of a finite set of (vertex) rays,
while a dually finitely generated cone is an $H$-cone, as it can be understood
as the intersection of finitely many linear halfspaces.

The repository already contained some versions of the Minkowski--Weyl theorem
in these terms, stated over a field.

%\begin{proposition}
A finitely generated cone $C$ over a field $\mathbb{k}$ can be rewritten
as the intersection $D \cap S$ where $D$ is a dually finitely generated
cone and $S$ is a finitely generated (linear) submodule.
%\end{proposition}

This holds even in infinite dimensional space because the result is proven
within the finitely generated submodule.
\begin{lstlisting}
variable (p) [Fact p.SeparatingRight] in
theorem FG.exists_dualfg_inf_submodule
  {C : PointedCone 𝕜 N} (hC : C.FG) {S : Submodule 𝕜 N}
  (hS : S.FG) (hCS : C ≤ S) : ∃ D : PointedCone 𝕜 N, D.DualFG p ∧ D ⊓ S = C := sorry
\end{lstlisting}

The hypothesis that $p$ is right-separating is there

Assuming finite dimensions, the result has a simpler statement.
\begin{lstlisting}
-- Assume N is finite dimensional
variable [Module.Finite 𝕜 N]

variable (p) [Fact p.SeparatingRight] in
lemma FG.dualfg {C : PointedCone 𝕜 N} (hC : C.FG) : C.DualFG p := sorry
lemma DualFG.fg {C : PointedCone 𝕜 N} (hC : C.DualFG p) : C.FG := sorry
\end{lstlisting}

\subsection{Definitions}
\label{subsec:mw:definitions}

% H-polyhedron
A \emph{polyhedron} is the finite intersection of halfspaces.
Classically, polyhedra are defined in $\mathbb{R}^n$ using a matrix inequality.
\[
  P = \left\{ x \in \mathbb{R}^n \;\middle|\; A x \le b \right\},
  \qquad A \in \mathbb{R}^{m \times n},\ b \in \mathbb{R}^m,
\]

Mathlib aims to avoid reliance on vector bases in statements and definitions,
given that they are not canonical, so instead our definition uses halfspaces
directly.

\begin{lstlisting}
variable {R : Type*} [CommRing R] [PartialOrder R] [IsOrderedRing R]
variable {V : Type*} [AddCommGroup V] [Module R V]
variable {V' : Type*} [AddCommGroup V'] [Module R V']
variable {A : Type*} [AddTorsor V A]

variable (R) in
def IsHPolyhedron (P : Set A) : Prop :=
    ∃ H : Finset (A →ᵃ[R] R), ∃ S : AffineSubspace R A,
      P = (⋂ h ∈ H, h ⁻¹' Set.Ici (0 : R)) ∩ S
\end{lstlisting}

Similarly to the existing \texttt{IsPolytope} definition, we use an existential
predicate on sets to define polyhedra.

% H-polyhedral cone
\begin{lstlisting}
def IsHPolyhedral
    (p : V' →ₗ[R] V →ₗ[R] R) (P : PointedCone R V) : Prop :=
  ∃C : PointedCone R V, ∃S : Submodule R V,
    C.DualFG p ∧ P = C ⊓ S
\end{lstlisting}

% Recession cone
Given a convex set $P$, its \emph{recession cone} is the cone of unbounded
directions of $P$.
For general sets, this definition is usually required in its strictest form,
that the recession cone added to each point of $P$ is contained in $P$.
\[
  \text{recess}(P) := \{ v \in V \mid
    \forall p \in P, \forall \lambda \in \mathbb{R}_+,
      p + \lambda v \in P \}
\]

Our formalization of this definition is straightforward.

% -- Could be for `ConvexSet` (to avoid polluting the `Set` namespace)
% -- This definiton assigns ⊤ as the recession cone of ∅
% -- The ∀ could be an ∃ (we could prove it's equivalent for polytopes)
% -- For `ConvexSet`, the ∀a is not necessary
\begin{lstlisting}
variable {k V W A : Type*}
variable [Field k] [LinearOrder k] [IsOrderedRing k]
variable [AddCommGroup V] [Module k V]
variable [AddCommGroup W] [Module k W]
variable [AddTorsor V A]

variable (k) in
def Convex.Set.recessionCone (P : Set A) : PointedCone k V where
  carrier := {v : V | ∀x ∈ P, ∀a : 𝕜, 0 ≤ a → a • v +ᵥ x ∈ P }
  sorry -- proofs that the carrier is a cone
\end{lstlisting}

We have defined the recession cone only for fields with a linear order.

\subsection{Minkowski--Weyl}
\label{subsec:mw:minkowski-weyl}

% TODO